\section{Notation and global inputs}\label{sec:setup}

All spaces are over $\C$ unless a reduction to positive characteristic
is specified. An irreducible holomorphic symplectic manifold is
abbreviated IHS. Write $q=q_X$ for the Beauville--Bogomolov--Fujiki
quadratic form, and $(\,\cdot\,,\,\cdot\,)$ for its associated
bilinear form. We also write $q(x,y)=(x,y)$, so that $q(x)=q(x,x)$.
We use the integral normalization; its real signature is
$(3,b_2(X)-3)$.

Choose a very ample line bundle $H$ on $B$, and set
\[
                         e=c_1(f^*H).
\]
For a coherent sheaf $M$ on a normal space, $M^{**}$ is its
reflexive hull. In particular,
$\Omega_B^{[j]}=(\bigwedge^j\Omega_B^1)^{**}$.
Reflexive tensor operations are always explicitly interpreted by
taking this hull. An associated map of a semiample line bundle is
replaced by its Stein factorization.

\begin{proposition}[Background for the given fibration]
\label{prop:background}
In the setting of Theorem~\ref{thm:main}, the following hold.
\begin{enumerate}
\item $X$ is projective, $e\neq0$ is nef and $q(e)=0$.
The base is $\Q$-factorial and klt, has Picard number one,
and $-K_B$ is ample as a $\Q$-Cartier divisor.
A smooth fiber is an abelian variety after choosing an origin.
\item There are reflexive identifications
\[
                  R^j f_*\cO_X\simeq\Omega_B^{[j]}.
\]
For every integral Weil divisor $A$ ample as a $\Q$-divisor,
\[
 H^i\!\left(B,
       (\Omega_B^{[j]}\otimes\cO_B(A))^{**}\right)=0
                       \quad(i>0,\ j>0).
\]
The reflexivity assertion also applies locally to the projective
equidimensional fibrations with connected fibers, smooth source \(Z\)
satisfying \(\omega_Z\simeq\cO_Z\), and normal Cohen--Macaulay
\(\Q\)-Gorenstein target that occur below.
\item The rational intersection cohomology of $B$ has the Betti
numbers of $\PP^n$, and the restriction of $H^2(X,\Q)$ to a
smooth fiber has one-dimensional image.
\item If $B$ is smooth, then $B\simeq\PP^n$.
\end{enumerate}
\end{proposition}

\begin{proof}
Relative ampleness and projectivity of $B$ make $X$ Moishezon;
as $X$ is Kähler, it is projective. Equivalently, a sufficiently
positive twist of a relative polarization by a base polarization
is ample. The structure statements in (1) are the fibration
theorems of Matsushita
\cite{Matsushita1999,Matsushita2001}.
The pullback of $H$ is nef and nonzero. Its $(2n)$-th power is
zero because $\dim B=n$, so Fujiki's relation gives $q(e)=0$.

For (2), Ou's Proposition~3.6 gives the reflexive cotangent
identifications for the given projective equidimensional Lagrangian
fibration, extending the symplectic differential identification of
Matsushita \cite{Matsushita2005,Ou2019}. Ou's Theorem~1.3 is the separate
local reflexivity result, and Theorem~1.4 gives the vanishing
\cite{Ou2019}.
The target hypotheses hold because klt spaces are
Cohen--Macaulay \cite[Theorem~11.1(2), Corollaries~11.9 and~11.14]{Kollar1997}
and $K_B$ is $\Q$-Cartier.
The cited reflexivity theorem has arbitrary-dimensional scope.
The higher-cohomology and fiber-restriction statements in (3)
are \cite[Theorem~0.4]{ShenYin2022}.
Statement (4) is Hwang's theorem \cite[Theorem~1.2]{Hwang2008}.
\end{proof}

Fix once and for all degreewise global isomorphisms
\[
 \phi_j:R^j f_*\cO_X\xrightarrow{\ \sim\ }\Omega_B^{[j]}
                  \qquad(0\leq j\leq n),
\]
with \(\phi_0\) the unit identification \(f_*\cO_X=\cO_B\).
The later normalized-chart constructions use the restrictions of these
same maps to \(B_{\mathrm{reg}}\); Proposition~\ref{str:direct-images}
constructs their compatible reflexive extensions on the charts.

We use intersection cohomology in Proposition~\ref{prop:background}
until rational smoothness has been proved in
Proposition~\ref{loc:rational-smoothness}, which then identifies it with
ordinary rational cohomology.

\begin{corollary}\label{cor:tangent-rigidity}
If $n>1$, then $H^1(B,T_B)=0$.
\end{corollary}
\begin{proof}
On the smooth locus, contraction identifies
$T_B$ with $\Omega_B^{n-1}\otimes\cO_B(-K_B)$.
Both sides extend reflexively, so
\[
 T_B\simeq
  (\Omega_B^{[n-1]}\otimes\cO_B(-K_B))^{**}.
\]
Apply Proposition~\ref{prop:background}(2), since $n-1>0$.
\end{proof}

The following elementary observation also appears in
\cite[Lemma~3.1]{GrebLehn2014}.

\begin{lemma}\label{lem:betti-bound}
One has $b_2(X)\geq4$.
\end{lemma}
\begin{proof}
Let $h$ be an ample class. Then $(h,e)>0$ and $q(e)=0$, so
$\langle h,e\rangle_{\R}$ is a hyperbolic plane.
It is perpendicular to the positive real plane generated by
$\operatorname{Re}\sigma$ and $\operatorname{Im}\sigma$.
These four independent real directions give the assertion.
\end{proof}

For a compact Kähler manifold $Y$, write $\mathcal K_Y$ for its
Kähler cone in $H^{1,1}(Y,\R)$.

\begin{theorem}[Strong hyperkähler SYZ, companion input]
\label{thm:strong-syz}
Let $Y$ be a compact IHS Kähler manifold of dimension $2m$.
If a holomorphic line bundle $L$ satisfies
\[
 c_1(L)\neq0,\qquad
 c_1(L)\in\overline{\mathcal K_Y},\qquad
 q_Y(c_1(L))=0,
\]
then a positive tensor power of $L$ is generated by its global
sections. Its associated map with connected fibers has a normal
projective base of dimension $m$. Every component of every fiber
has dimension $m$ and is Lagrangian.
\end{theorem}

This is Theorem~1.1 of the companion \cite{StrongSYZ2026}; it has no
restriction on \(b_2(Y)\), \(m\), or deformation type. The hypothesis
\(b_2\geq5\) in that manuscript belongs to its separate finiteness
Corollary~1.2. Section~\ref{sec:second-system} applies the theorem to a
second nonzero integral nef isotropic class on an actual projective
deformation. The class is of type \((1,1)\), so the exponential sequence
supplies its holomorphic line bundle.

\paragraph{Reduction to smoothness.}
For $n=1$, normality makes $B$ a smooth curve, and
Proposition~\ref{prop:background}(4) proves the result.
Henceforth $n>1$. It suffices to prove that $B$ is smooth.
Proposition~\ref{def:fixed-base} later supplies proper holomorphic
fibrations over this same \(B\) throughout the nearby \(e\)-Hodge locus.
The original fibration and the projective deformation carrying two
Lagrangian systems have the stronger projectivity used in their respective
arguments.

\paragraph{The Euler obstruction.}
The first stage has a specific numerical target. If \(B\) has a
nonquotient germ, Sections~\ref{sec:local-tensor}--\ref{sec:stratum}
produce a nonempty closed irreducible normal projective stratum
\(S_0\subset B\) of dimension \(d<n\). Proposition~\ref{eu:euler-zero},
proved through Sections~\ref{sec:base-change}--\ref{sec:euler}, establishes
\[
                         \chi_{\mathrm{top}}(S_0)=0.
\]
Here \(\chi_{\mathrm{top}}\) is the ordinary topological Euler
characteristic of the coarse space. The local construction also gives a
smooth stack \(\cS\) above \(S_0\); we write \(c=n-d>0\).
Its finite stabilizers, including those acting trivially on the stratum,
are retained in the calculations. In characteristic \(p\), their orders
are chosen prime to \(p\). The later geometric arguments will contradict
this Euler equality when \(b_2(X)\geq5\) and in the isotrivial and varying
cases with \(b_2(X)=4\).
